\section{Expectation-to-failure-probability via Markov}
\label{sec:expectation_to_failure}

\begin{lemma}[Markov bound on the failure probability]\label{lem:expectation_to_failure}
Let $X$ be any random variable taking values in $\R^d$ (e.g.\
$X = x_T$ for a fixed terminal index, or $X = x_R$ for a randomised
iterate). For every $Q \in F$,
\begin{align}\label{eq:markov_loss_tail}
   \Pr\!\bigl[\, \loss(X; Q) \;>\; \log 2 \,\bigr]
   \;\le\;
   \frac{\E[\loss(X; Q)]}{\log 2}.
\end{align}
Consequently, on the event that the non-degeneracy condition
\Cref{eq:logit_margin_nondegen} of \Cref{lem:logit_margin_decoding}
holds for $X$,
\begin{align}\label{eq:markov_decoding_failure}
   \Pr\!\bigl[\, \dec(X) \notin \Cset(Q) \,\bigr]
   \;\le\;
   \Pr\!\bigl[\, \loss(X; Q) \;>\; \log 2 \,\bigr]
   \;\le\;
   \frac{\E[\loss(X; Q)]}{\log 2}.
\end{align}
\end{lemma}


\begin{proof}
We prove the two displays in order.

\noindent\textbf{Step 1: Markov on $\loss(X; Q)$.}
By \Cref{rem:loss_nondegenerate}, $\loss(X; Q) \ge 0$ almost surely (the
non-negativity is uniform in $x$ and hence transfers to any random $X$).
Markov's inequality applied to the non-negative random variable
$\loss(X; Q)$ at threshold $\log 2 > 0$ gives
\begin{align*}
   \Pr\!\bigl[\, \loss(X; Q) > \log 2 \,\bigr]
   \;\le\;
   \frac{\E[\loss(X; Q)]}{\log 2},
\end{align*}
which is \Cref{eq:markov_loss_tail}.

\noindent\textbf{Step 2: bound on the decoding-failure event.}
On the event $\{\loss(X; Q) \le \log 2\}$, both implications of
\Cref{lem:logit_margin_decoding} apply (the second using the assumed
non-degeneracy condition \Cref{eq:logit_margin_nondegen} of
$X$), so $\dec(X) \in \Cset(Q)$. Contrapositively, the decoding-failure
event $\{\dec(X) \notin \Cset(Q)\}$ is contained in
$\{\loss(X; Q) > \log 2\}$. Therefore
\begin{align*}
   \Pr\!\bigl[\, \dec(X) \notin \Cset(Q) \,\bigr]
   \;\le\;
   \Pr\!\bigl[\, \loss(X; Q) > \log 2 \,\bigr],
\end{align*}
and combining with Step 1 yields \Cref{eq:markov_decoding_failure}.
\end{proof}

\begin{remark}\label{rem:markov_tightness}
Markov's inequality is the simplest tail bound and is generically loose
by a constant factor (Chebyshev's inequality with a second-moment bound
on $\loss(X; Q)$ would tighten the $1/\log 2$ to a $1/(\log 2)^2$ factor
against a second-moment quantity). For the present purpose we adopt
Markov because (i) it requires only \Cref{ass:initial_loss} and the
descent analysis of \Cref{lem:telescoping}, no second-moment hypothesis
on $\loss(X; Q)$, and (ii) the resulting failure-probability bound
\Cref{eq:markov_decoding_failure} has the same asymptotic rate
$O(1/\sqrt T) + O(\beta^2/\eta_0)$ as the loss bound, so improving the
constant in front does not change the qualitative picture.
\end{remark}
